% Chapter VII: Drawing Methods in Byzantine and Old Russian Painting
% Source pages: 140–167

\chapter{Drawing Methods in Byzantine and Old Russian Painting}

\markboth{Chapter VII. Drawing Methods}{Chapter VII}

\srcpage{140}

The art of ancient Egypt, examined in Chapter~I, was based on reproducing the geometry
of objective space.  Medieval art turned to conveying the geometry of visible, perceptual
space.  Unlike the painting of the Renaissance, which pursued similar goals while at the
same time strictly excluding any recourse to the geometry of objective space, medieval
painting admitted such ``liberties.''

As recalled from Chapter~I, conveying the geometry of objective space requires drawing
methods, and accordingly in medieval art one encounters examples of a symbiosis of
technical drawing and pictorial rendering.  Making no fundamental distinction between the
two, the medieval master skillfully combined these two different approaches in his pursuit
of the best solution for artistic and informational expressiveness.

Among the images created by Old Russian painters one can identify works in which the
primary task was to convey not the visible but the \emph{objective} geometry of space.
First let us consider those that today we would call topographic plans.  An icon of the
mid-seventeenth century, for example, depicts the Trinity-Sergius Monastery---clearly a
plan of the monastery, in which the artist evidently aimed to convey the true disposition
of walls, towers, churches, and all other buildings, omitting not a single one, however
secondary (Ill.~59).  Yet this plan is far from a modern one.  Striking immediately is the
fact that all buildings are shown in a conventional rotation, viewed from the side, as a
person arriving at the monastery would see them.  This device is entirely natural; it was
widely used in ancient Egyptian painting and, as already noted, lives on today in sketch-maps
published for tourists.  The drawing method in this icon is so accomplished that an ancient
Egyptian would unquestionably have been able to ``read'' it with ease.  Even so, the image
is not free from elements of pictorial rendering, though these play a purely subordinate
role---visible, for instance, in the foreshortened earth at the bottom of the icon, in
the rendering (contrary to drawing rules) of the lateral walls of the Assumption Cathedral,
and in the clear perspectival display of its five domes.

\figblock{59}{``The Trinity Monastery.'' Icon of the mid-seventeenth century.
  Zagorsk Historical-Artistic Museum Reserve (ZIKHM).}

\srcpage{141}

\figblock{60}{Foundation of the stone church of St~Michael in Novgorod.
  Miniature from the \emph{Drevny letopisets} [Ancient Chronicle], vol.~II,
  fol.~494, sixteenth century.}

Consider an example of the opposite type.  A miniature of the second half of the
sixteenth century depicting the laying of the foundations of a church---which, by the
definition introduced earlier, is a picture---contains a plan of the foundations of the
church being erected.  The introduction of drawing methods into the picture allowed the
miniaturist to convey what was, from his point of view, important information: the reader
of the chronicle was told that a three-apse, four-pier stone church was being founded
(Ill.~60).

In the examples examined one sees how drawing and pictorial methods, applied jointly,
enhance the informational content of both drawing and picture alike.  The main concern
of what follows, however, is the narrower question of drawing methods that complement
the picture.

The typical drawing devices that find application in medieval painting are:
(1)~conventional rotations of projection planes, (2)~cross-sections, and (3)~unfoldings.
We shall examine these three drawing methods in the order named.

\srcpage{142}

\noindent\textbf{1.}\quad\textbf{Rotations of projection planes.}
The conventional rotations in question have already been discussed in Chapter~I.  Here
we describe certain features of this method of depiction that are characteristic of Old
Russian and Byzantine art but not found in ancient Egyptian art.  As already noted in
Chapter~I, for an object of complex configuration one may depict individual parts of it
in conventionally rotated positions.  In essence, this officially sanctioned drawing
method existed in medieval art when---motivated by the desire to convey more
information---certain planes of an object were depicted in conventionally rotated positions
in the picture as well.  Such conventional rotations were discussed in the preceding
chapter and need not be described in more detail here.

Another variant for depicting on a drawing an object whose configuration cannot be
understood from a single projection is the combined depiction of two of its projections
together---that is, showing it ``from the side'' and ``from above'' simultaneously.
Here too, in effect, one is dealing with a conventional rotation of one of its planes.
Consider the engineering drawing of a section of circular pipe with a flange (shown with
a minor deviation from current standards) (Ill.~61).  The lower part of the drawing shows
the side view, which, though necessary, does not allow the shape of the flange or its
other features (bolt-holes) to be established.  These are visible in the upper part of
the drawing, where, by reason of symmetry, only half the flange is shown.  Such a
representation is compact, simple, and clear.  One might even maintain that the drawing
shown is virtually the most natural way of depicting objects of this kind.

\figblock{61}{Engineering drawing of a flanged pipe section.}

\figblock{62}{Table in the fresco ``Trinity'' from the Kiev Cathedral of St~Sophia
  (traced drawing).}

\srcpage{143}

This last claim is supported by the fact that essentially the same method of depiction
is found quite frequently in Old Russian art.  In the frescoes of St~Sophia Cathedral
in Kiev (eleventh century), ``The Miracle at Cana of Galilee'' and
``The Trinity,''\footnote{G.~N.~Logvin.  \emph{Sofia Kievskaya} [Kiev Sophia].
  Kiev, 1971, Ills.~226, 244.} the front edges of the round tables are shown as
horizontal lines with tablecloths hanging from them---that is, viewed from the side---while
their horizontal surfaces are shown as semicircles, viewed from above, following a scheme
corresponding to drawings of the type shown in Ill.~61.  Illustration~62 reproduces a
drawing of the table after the ``Trinity'' fresco, conveying the main geometric features
of its representation.\footnote{A natural question arises here: might not semicircular
  tables have actually existed?  Such tables did indeed exist in antiquity (when people
  reclined rather than sat at table); they are sometimes called sigma-shaped.  However,
  they disappear entirely from use at the very beginning of the Middle Ages, and medieval
  artists, rendering a table top as a semicircle, clearly understood that they were
  depicting a round table.  This is confirmed, among other things, by the fact that round
  footrests are likewise rendered as semicircles, even though sigma-shaped footrests never
  existed.}  In the panel depicting the ``Wedding at Cana of Galilee'' on the Novgorod
icon ``The Earthly Life of Christ,'' the round table is also shown in two projections
simultaneously---in side view (straight outlines) and in top view (the curved outline of
the tabletop) (Ill.~63).  The side view is rendered very clearly, with vessels standing
on the table placed on the straight-line part of the table's outline, since that is
exactly how they appear when the table is viewed from the side.  As for the horizontal
surface of the table, though shown in top view, the direction of gaze onto it is not
perpendicular to the side-view direction (as in Ill.~62) but deviated from it by
approximately $30°$ from the horizontal in this case;

\srcpage{144}

this is perfectly natural for later works, since introducing foreshortening, though moving
the image away from the objective geometry of space, brought it closer to the picture---to
the visible geometry of space.  In the examples considered, the artist's desire to convey
to the viewer the characteristic features of the table from different directions of gaze
is clearly evident.

One should not suppose that this type of depiction of a round table was rooted in the
artist's inability.  In the festal tiers of the iconostases of the Annunciation Cathedral
in Moscow and the Trinity Cathedral of the Trinity-Sergius Lavra, created with the
participation of Theophanes the Greek, Prokhor of Gorodets, Andrei Rublev, and Daniil
Chyorny, there are depictions of tables on the icons of the ``Last Supper'' shown as
complete ovals, roughly as a modern artist would draw them.\footnote{The depiction of a
  round table as an oval was no novelty for Russian artists of the fifteenth century.
  This type of representation appears, for example, in the mural painting of the
  Saviour-Transfiguration Cathedral of Mirozhsky Monastery (twelfth century, Pskov) in
  the composition ``Conversation of Joachim and Anna.''  See: M.~N.~Soboleva.
  ``Stenopis' Spaso-Preobrazhenskogo sobora Mirozhskogo monastyrya v Pskove.''
  \emph{Drevnerusskoe iskusstvo}.  Moscow, 1968, p.~20.}  And yet Theophanes the Greek
painted his ``Trinity'' in the Church of the Transfiguration of the Saviour in Novgorod
with the figures seated at a conventionally-depicted, drawing-method table, and in the
frescoes of the Church of the Nativity of the Virgin in Ferapontov Monastery, created by
Dionysius approximately eighty years after the completion of work on the Trinity Cathedral
iconostasis, one can see numerous examples of the simultaneous depiction of an object
``from the side'' and ``from above,'' in the manner of the tables on the Kiev frescoes
and the Novgorod icon described above.  In Dionysius's fresco ``Wedding at Cana,'' not
only the table but also the round footrest of Christ is shown by the method described,
with the contours of the horizontal surface of the footrest closer to a semicircle than
those of the table---it appears to be seen from a steeper direction, at approximately
$45°$ rather than the $\sim\!30°$ of the table (Ill.~64).  While the horizontal surfaces
of the tables in the other frescoes of this church are always shown foreshortened at
$\sim\!30°$, the same cannot be said of the round footrests, which for understandable
reasons are always depicted from a steeper angle and sometimes almost exactly from above,
that is, as almost exact semicircles.\footnote{See: I.~E.~Danilova.  \emph{Freski
  Ferapontova monastyrya} [Frescoes of Ferapontov Monastery].  Moscow, 1969,
  Ills.~133, 54.}  As examples one may cite the Ferapontov frescoes ``The Fourth
Ecumenical Council'' and ``The Parable of the Wise and Foolish Virgins,'' where the
horizontal surfaces of the round footrests are rendered as if viewed exactly from above,
while their legs are shown in side view.

\figblock{63}{``The Wedding at Cana of Galilee.''  Panel of the icon ``The Earthly Life
  of Christ.''  Fifteenth century.  Novgorod State Museum (NGM).}

\srcpage{145}

Sometimes it seemed more natural to the artist to render the front edge of the table
not as a straight line but with a slight curvature, as if introducing a mild foreshortening
into the side view as well.  This is done in certain frescoes of Dionysius, for example
in the fresco ``The Feast at Simon the Leper's.''\footnote{I.~E.~Danilova.
  \emph{Op.~cit.}, ill.~89.}

The conventionally-drawn depiction of round tables occurs in Old Russian painting far
more often than their natural rendering as ovals.  This is apparently connected with the
fact that this type of depiction brought the figures seated at the table closer to the
viewer.  For if the table in the Novgorod icon ``Wedding at Cana of Galilee'' (Ill.~63)
were rendered as an oval, Christ would appear in the depths of the depicted space rather
than in the immediate foreground.  The same considerations apply to the Dionysius fresco
cited, where not only the table but also the round footrest of Christ is rendered so as
not to distance Christ from the viewer.

\figblock{65}{Depictions of baptismal fonts.  Upper: from the Novgorod icon ``Nativity
  of the Virgin''; lower: from the fresco ``Healing of the Blind Man'' in the Church of
  the Ascension, Ravanica (Yugoslavia).  Schematic renderings.}

The method of simultaneously showing ``the side view'' and ``the top view'' was not
confined to tables and footrests.  On the Novgorod icon ``Nativity of the Virgin'' the
baptismal font for the washing of the newborn infant is depicted in the same way.  An
analogous depiction of a font appears in the fresco ``Healing of the Blind Man'' in the
Church of the Ascension, Ravanica, Yugoslavia.  Both fonts are shown schematically in
Ill.~65.

\srcpage{146}

\figblock{64}{Dionysius.  ``The Wedding at Cana of Galilee.''  Fresco of Ferapontov
  Monastery, sixteenth century (traced drawing).}

Examples of depicting an object, or a part of it, in a conventionally rotated position
could be multiplied without difficulty, so widespread are they.  Let us note here only
the depictions of vessels on tables (Ills.~52, 53, 62).  These vessels are shown in side
view (clearly visible from their straight-line bases), except for the upper parts, which
are depicted in foreshortening.\footnote{This is well seen in the panels of the icon
  ``Nicholas in His Life'' from Ozerevo, held in the State Russian Museum.  Noteworthy
  is that the cups placed on the tables are not only depicted exactly on the front edge
  but are painted in the same colour as the ``lateral projection'' of the table, thereby
  emphasising their belonging to that projection.  The device of painting vessels shifted
  to the straight front edge of the table in the same colour as the lateral projection
  is fairly widespread; see, for example, the icon ``Old Testament Trinity'' from Rostov
  the Great (N.~V.~Rozanova.  \emph{Rostovo-suzdalskaya zhivopis XII--XVI vekov}
  [Rostov-Suzdal Painting of the Twelfth to Sixteenth Centuries].  Moscow, 1970,
  ill.~28).}  This method of depiction is not without clarity: since all the vessels
shown are axially symmetric, they are adequately characterised by their side view, while
the oval upper parts convey the functionally important features of the vessels, such as
the openings in the necks of bottles, and so on.

\srcpage{147}

The predominant depiction of vessels in side view sometimes leads---especially in earlier
works---to their being placed on the front edge of the table (Ill.~63), that is,
incorporated into the common side view with the table.\footnote{G.~N.~Logvin.
  \emph{Op.~cit.}, Ills.~244, 152.}  This shift of vessels produces the characteristic
``displacement toward the viewer'' of depicted objects, already described in the
literature on medieval painting.  Even in cases where vessels standing on tables (or
altars) are displaced in depth relative to the front edge, the tendency to depict them
close to this edge persists.  Here the tradition of placing vessels so that their
silhouette lies entirely within the table's outline may also play a role.

Exceptions are occasionally observed only in later works and for vessels with very tall
necks.  This probably seemed entirely logical to medieval people.  At any rate, in the
mural painting of St~Sophia Cathedral in Kiev, the cups standing at the front edge of
the table in the fresco ``Trinity'' are sized so that their upper rims do not pass beyond
the depiction of the far boundary of the table as seen by the viewer (see Ill.~62).  As
a result the middle cup is much larger than the lateral ones.  The distribution of objects
on tables is always such that large pieces (cups, goblets) are displaced toward the front
edge, while only small ones are placed at the back (see Ills.~27, 52, 53).  The artist
often pays little attention to ``table-setting'' as such and simply ``lists'' the objects.
In the Kiev Sophia fresco ``Peter's Baptism of the Daughter of the Centurion Cornelius''
the contours of the font are distorted; its far edge is displaced so far from the viewer
that it cannot be explained by rendering the font according to the scheme of Ill.~65.
Clearly it was important that the figure of Cornelius's daughter should not extend beyond
the silhouette of the font.  An analogous phenomenon can be observed in the miniature
``Baptism'' from a manuscript in the Matenadaran (Yerevan; No.~7639, fol.~5a).  Here
Christ and John stand in the river that flows along the lower edge of the miniature, and
therefore Christ's figure ought to be shown against the bank.  In fact the river seems
to widen at the figure of Christ and he finds himself in a ``cove,'' in complete
conformity with the dimensions of his figure.  Essentially analogous is the composition
of the ``Baptism'' in Ill.~35.

Although the considerations presented here bear no direct relation to methods of spatial
construction, they do influence the composition, which in turn---as already noted---bears
on spatial construction.

\srcpage{148}

\noindent\textbf{2.}\quad\textbf{Cross-sections.}
In order to increase the informational content of a representation, various kinds of
cross-sections and cuts can be employed.  The effectiveness of this method was already
demonstrated in the discussion of the features of ancient Egyptian painting.  It is
employed today in technical drawing as well, when a drawing of the exterior is supplemented
by a representation of the internal structure through the conventional removal of part of
the construction and the display of the features concealed behind the removed portion.
Let us cite an example of such a drawing using what is called a cross-section (Ill.~66).
It shows a hypothetical device containing inside, on a circuit board~$B$, an electronic
circuit.  This board is housed in a casing with front cover~$A$.  To make the internal
structure of the unit visible, the front cover is conventionally removed between the
lines~$C$, which are the boundaries: on one side of them the external cover of the device
is shown, and on the other the board inside it is depicted.

\figblock{66}{Engineering drawing using a cross-section: a hypothetical electronic
  device.  The front cover~$A$ is conventionally removed between lines~$C$, revealing
  the interior circuit board~$B$.}

The possibility of showing what in real life cannot be seen simultaneously but
nonetheless exists simultaneously---this possibility was bound to attract medieval artists
who were not bound by the requirement of an illusionistic rendering of the external world.

On the Russian icon ``Deposition of the Holy Robe'' of the seventeenth century (Ill.~67),
the Assumption Cathedral of the Moscow Kremlin is depicted: its exterior is shown in the
upper part, and the baldachin inside the cathedral and the people within it are shown in
the lower part.  This was possible only because the artist removed part of the outer
wall---doing precisely what the draughtsman does in Ill.~66.  This device, with various
minor variations, became very widespread in Russia in the seventeenth century, as evidenced
by the numerous panels of hagiographic icons and the frescoes of the churches of Rostov,
Yaroslavl, Kostroma, and in Moscow, for example, the painting of the Church of the
Trinity in Nikitniki.

Subjects connected with religious legend set before medieval masters the task of
simultaneously depicting on a single picture events occurring in the real three-dimensional
world and in the mystical world---the latter conceived as also three-dimensional but
immaterial, though in the Middle Ages few doubted the existence of this invisible world.
Moreover, both these three-dimensional worlds, the ``visible'' and the ``invisible,'' were
often supposed to exist not only simultaneously but also in one and the same region of
space, and not independently of each other but in close interaction.

\srcpage{149}

To give a mathematical interpretation of what has just been said, one may maintain that
the artists were required to depict a \emph{four-dimensional space}.\footnote{Four-dimensional
  space is often understood as a space with three ordinary coordinates and time as the
  fourth.  Here we are considering the case where there are four spatial coordinates,
  and time (if it too is interpreted as a coordinate) would be the fifth.

  One should guard against a simplified understanding of multi-dimensionality.  One
  cannot reason as follows: a picture depicts a three-dimensional world, and before the
  picture stands the viewer in another, also three-dimensional world, so altogether the
  system ``viewer and picture'' has six dimensions.  Each additional dimension changes
  not the quantitative but the qualitative character of space.  A line is a one-dimensional
  object, but two lines do not yield a two-dimensional space---that is a qualitatively new
  entity, a plane.  The transition to three dimensions gives a new quality---volume; and
  so the recourse to four-dimensional space must give the next, fourth level in the
  qualitatively escalating sequence: line, plane, volume\ldots}  Modern geometry knows
multidimensional spaces and studies their properties.  However, although the formal
properties of such spaces are susceptible of description, depicting them is impossible,
because the real space known and observed in ordinary human experience is three-dimensional.
The most direct way of feeling the insuperable difficulties that arise when one attempts
to visualise four-dimensionality is to take three mutually perpendicular straight lines
issuing from a single point---the ordinary coordinate axes of space---and attempt to
construct a straight line also issuing from the origin but simultaneously perpendicular
to all three of these ordinary axes.

If one were to give the chief characteristic of four-dimensional geometric space---chief
from the standpoint of analysing medieval images---it amounts to the assertion that within
our ordinary three-dimensional volume, for instance the volume of a room, different
``lives'' can exist independently and simultaneously.

Let us agree that the world we are considering is four-dimensional.  Then three coordinates
determining the position of some point in the space of a room (distances in three mutually
perpendicular directions---forward, sideways, and upward---from some starting point) are
insufficient to fully determine its position, since by the adopted assumption of
four-dimensionality one must also specify its distance in the ``fourth'' direction
(likewise ``perpendicular'' to those named above).  It is impossible to visualise this
``fourth'' direction, which takes the point out of the three-dimensional space familiar
to us.  The impossibility of visualisation does not at all mean the impossibility of
analysing the properties of such a space or of describing them clearly.

Let each point of the ordinary space of a room correspond to a coordinate~$w_1$ measured
along the ``fourth'' axis.  Now take, in this ``fourth'' direction, points with
coordinate~$w_2$.  The result is that each point of our ordinary space must be counted
twice: once with fourth coordinate~$w_1$, and again with~$w_2$.  Clearly, these two
different regions of four-dimensional space (let us call them space~$w_1$ and
space~$w_2$) will not coincide; they will be located in different parts of the
four-dimensional space, and different events may occur in space~$w_1$ and in
space~$w_2$.

\srcpage{150}

\figblock{67}{``Deposition of the Holy Robe.''  Icon of the 1630s.
  State Historical Museum (GIM).}

As for the distance between spaces~$w_1$ and~$w_2$ along coordinate~$w$, it may be
arbitrarily small, and yet the two spaces will not ``merge.''  This is impossible to
visualise directly, so let us explain by analogy: if one takes an infinitely thin plane
(for instance a sheet of paper), different events may occur on its two sides without
interfering with each other, even though these two planar spaces are displaced relative
to each other by an infinitely small amount---the ``thickness'' of the plane in
question.\footnote{Those who wish to acquaint themselves more thoroughly with the
  properties of multidimensional spaces are referred to the numerous popular-science works
  of mathematicians and philosophers.

  A brief and entirely elementary treatment of the geometry of four-dimensional space,
  sufficient for understanding the questions raised here, is given in Appendix~10.}

Consequently, if one introduces an infinitely small displacement along coordinate~$w$,
then to each point of the three-dimensional space of the room there will correspond two
infinitely close but non-coincident regions of four-dimensional space.

The geometric constructions described above formally possess the properties needed to
describe many events that form the content of medieval painting.  Indeed, let space~$w_1$
be the space in which the people and other beings that come into direct contact with the
people depicted in the icon act, and let space~$w_2$ be the mystical space in which the
incorporeal inhabitants of the mystical realm---angels and so on---live and act.  In that
case, both bodily and incorporeal characters may act in one and the same room without
``getting in each other's way,'' since they are located in different regions of
four-dimensional space.  This ability to preserve the territorial unity of mystical and
real while at the same time granting each its necessary independence answers completely
to the medieval conception of the relation between the real, earthly and the mystical.

If one attempts to show these two spaces together visually, a method may be proposed that
reduces to the \emph{method of sections}, whose essence is the interruption of the
depiction of one space when it becomes necessary to turn to the depiction of the other.
The simplest case in terms of content merely reflects the principle of depicting two
worlds when they are clearly separated---divided---by the very subject of the composition.
Although these two worlds are geometrically separated, they interact.  The ideas connected
with this circumstance are best considered together with the general ideological and
figurative content of medieval painting, which lies beyond the scope of the present book.
We are concerned here only with the logic of the construction of spatial relations, and
we see how common sense and clarity of method sometimes stand in a strangely paradoxical
juxtaposition with the complex irrational meaning of the image.

\srcpage{151}

Let us turn to the Novgorod icon ``Dormition,'' attributed to Theophanes the Greek or in
any case painted under his strong influence (Ill.~68).  The masters who painted the icons
of the ``Dormition'' faced precisely the task discussed above.  According to church
tradition, the Mother of God died surrounded by grieving apostles, and her soul (depicted
as an infant) was taken to heaven by Christ, who appeared for this purpose.  The important
point is that these two events---the passing of Mary surrounded by the apostles standing
beside her bier, and the taking of her soul by Christ---occurred simultaneously, in one
space, but the first in the real plane and the second in the mystical.  From the standpoint
of the formal principles developed just above, these two events are conveniently treated
as occurring in one three-dimensional space but in different regions of the four-dimensional
space: the visible~$w_1$ and the invisible~$w_2$.

Using the method of sections, both regions can be depicted simultaneously by dividing the
picture plane into two parts, just as the drawing plane in Ill.~66 was divided for an
analogous purpose into regions~$A$ and~$B$.  On the drawing, regions~$A$ and~$B$ are
separated by the boundary~$C$; by analogy, on the icon in question, regions~$w_1$
and~$w_2$ must be separated from each other in a clear and unambiguous
manner.\footnote{Persons familiar with the terminology used in technical drawing might
  suppose that if the space in the ``Dormition'' icon is to be interpreted as a drawing
  at all, it would be more correct to speak of a ``cut'' rather than a ``section.''  But
  this is not so: in four-dimensional space, a section of it (the fixing of one
  coordinate) yields not a plane (as in ordinary three-dimensional space) but a volume.}

On the Novgorod icon this boundary is the line separating the ordinary light background
of the sky from the dark-blue background of the space in which Christ stands.  Clearly,
this dark-blue space has been singled out by the method of sections, since it interrupts
the depiction of the architectural background.\footnote{The simplification of
  distinguishing between the external and the sectionally revealed internal structure of
  an object by colouring their images in different colours is itself a drawing device.
  It was used as late as the end of the nineteenth and the beginning of the twentieth
  century for particularly important drawings.  At the present time, instead of colouring
  a section in a different colour, various types of conventional hatching are used.}  It
turns out that the icon depicts two spaces: the real, to which the bier of Mary, the
apostles, the prelates, and the architectural background belong; and the mystical,
containing Christ, which begins between the foreground and background of the real
space.\footnote{The terms ``real'' or ``ordinary'' space are of course conventional.
  As is well known, the space depicted in icons was not the everyday space of human
  practice; but qualifications of this kind would be superfluous here.}

\figblock{68}{Theophanes the Greek~(?).  ``Dormition.''  Icon of the late fourteenth
  century.  State Tretyakov Gallery (GTG).}

\srcpage{152}

The artist emphasises in every possible way that these two spaces are connected only
through the mystical act---the taking of Mary's soul---and that the mystical space remains
invisible to those gathered around Mary's bier.  The gaze of every figure is directed
toward the dying Mary, and no one looks at Christ, even though the appearance of Christ
among his disciples would obviously have overwhelmed them, the more so as he appears in
radiant and luminous garments, as he is depicted on icons of the ``Transfiguration.''

The desire to mark the boundary of the section as clearly as possible---so as
simultaneously to show the two coexisting regions of space---led to the fact that from
the very beginning of the fifteenth century this boundary is particularly often emphasised
not only by colour but by the depiction of an unbroken row of angels disposed along the
border between the two depicted spaces.  On the one hand these are the heavenly hosts
accompanying Christ; on the other, the fundamental difference between the two layers of
the four-dimensional space was underscored, the passage of the border between which was
barred by the heavenly hosts.  These angels are painted monochromatically against the
background of the mystical space, which symbolises their indissoluble connection with it
and simultaneously makes the angels as it were ``invisible,'' in keeping with the nature
of the mystical space they bound.

In the fifteenth and sixteenth centuries a preference arose for more complex compositions
for the ``Dormition'' icons; in particular, the region of mystical space is often shown
three times, in different areas of the picture plane.  Such is the icon ``Dormition'' from
the Belozersky Monastery at Kirillov, attributed to a master of Andrei Rublev's circle
(Ill.~69).  Here the section of the four-dimensional space for the purpose of showing
the mystical is made three times: first, at the depiction of Christ with ``invisible''
angels disposed along the boundary of the section; a second time at the depiction of Mary
being carried heavenward by angels; and a third time in the area of the upper edge of the
icon, in the form of a semicircle depicting the mystical sky---the threshold of
paradise---also with ``invisible'' angels.  All three of these regions, by their very
meaning, cannot belong to the main space of the icon, and to emphasise this the artist
uses for all of them related, fundamentally blue colours sharply distinct from the gold
background by which the sky---or more precisely, the heavenly light---is rendered.

The introduction of a colour for the mystical space distinct from the remaining parts of
the image, and its delineation by clearly drawn boundaries, rested on a clear logic.  In
earlier works the logic of the coexistence of two parallel spaces had not yet been so
strictly thought out.  For example, on the Novgorod icon ``Dormition'' from the Desyatinny
Monastery (twelfth to thirteenth century), Christ is depicted against the same gold
background as the apostles; only in the upper part of the icon does the dark mystical sky
of paradise, bounded by an arc of a circle, show blue.  Christ, slightly elevated above
the group of apostles, barely distinguishes himself from their company, and it remains
unclear why the apostles and prelates do not notice him.\footnote{The image of this icon
  is reproduced, for example, in the album \emph{Shedevry drevnerusskoy zhivopisi}
  [Masterpieces of Old Russian Painting] (Moscow, 1971, ill.~3).}  In the later icons of
the sixteenth and seventeenth centuries the composition of the icons sometimes grows more
complex while the logical clarity of the fourteenth- and fifteenth-century icons is lost.
By the seventeenth century there appear icons of the ``Dormition'' in which different
colours are assigned to the mystical space in different parts of one and the same image.
The artist is clearly drawn to ornate polychromy and has already forgotten the strict
meaning of the introduction of conventional colours for distinguishing ordinary from
mystical space.

\srcpage{153}

\figblock{69}{``Dormition.''  Icon of the Kirillo-Belozersky Monastery.
  Circle of Andrei Rublev.  First third of the fifteenth century.
  Pereyaslavl-Zalessky Museum (PT).}

In the works of Theophanes the Greek, Andrei Rublev, Dionysius, and the artists close
to them, logically impeccable compositions on the theme of the ``Dormition'' still today
evoke the impression of a mathematically rigorous construction by their inner freedom from
contradiction.  Without knowledge of modern multidimensional geometry, the great masters,
with their inherent artistic intuition and rigorous reasoning, justly earned the right to
be called ``most illustrious sages, philosophers exceedingly skilled''---even if one
confines oneself to the geometrical aspect alone.\footnote{A natural doubt may arise here:
  do not these arguments oblige us to attribute very deep mathematical knowledge to
  medieval artists?

  Of course, Russian artists of the fifteenth century did not practise mathematics;
  however, confronted with the problem of the non-contradictory coordination of mystical
  tradition and everyday human experience, they were compelled to construct a certain
  logical system, which turned out to be isomorphic to the geometry of four-dimensional
  space.

  It is not superfluous, finally, to note that even today the depth of medieval
  scientific thought, which, as is well known, wore a theological guise, is conventionally
  underestimated.  In this connection one may recall, for example, that as far back as
  the fifth century Blessed Augustine taught the simultaneous creation of the world and
  of time by God.  This conception ``made an important contribution to the philosophical
  theory of time and even in some respects anticipated Einstein's idea of the relativity
  of time,'' according to which time without matter does not exist (L.~M.~Mostapenko.
  \emph{Prostranstvo i vremya v makro-, mega- i mikromire} [Space and Time in the Macro-,
  Mega-, and Microworld].  Moscow, 1974, pp.~180--181).}

\srcpage{154}

The use of the method of sections in the simultaneous depiction of ordinary and mystical
space is of course not confined to compositions of the ``Dormition.''  In confirmation,
let us cite other subjects.  Very often on icons and frescoes of the ``Annunciation'' a
blue arc of a circle is shown at the upper edge of the image, symbolising the mystical
sky, from which blue rays extend toward Mary, symbolising her mystical connection with
this sky.  Exactly the same connection is depicted on many icons of the ``Nativity of
Christ,'' on icons of ``John and Prokhor on the Island of Patmos,'' and so on.

When individual saints are depicted, especially when shown at prayer, in the middle of
the upper edge of the icon or more often in one of its upper corners there is a segment
of a circle---usually painted dark blue---in which Christ, a blessing hand of the deity,
or something similar is depicted.  The icon of ``Anthony of Rome'' from the sixteenth
century gives an example of such a composition (Ill.~70).  In it Anthony raises prayers
to the Mother of God and to Christ, who are shown in the upper left corner of the icon,
with the sky of the space in which Anthony stands and the mystical sky in which the
Mother of God with the infant Christ is visible sharply distinct in colour.  In all these
cases the boundary between real and mystical space is drawn very clearly; both spaces
always differ sharply in colour; sometimes their boundary is emphasised by a row of stars
disposed along it on the dark background of the mystical sky, sometimes by one or more
coloured stripes, or in some other way.  In the cases described, we are of course dealing
with situations where the figure depicted on the icon does not see the deity and is
connected with the mystical space only in thought.

In cases where the deity or a saint enters into visual contact with people depicted on
the icon or fresco---that is, as it were enters the real world---the method of sectioning
the four-dimensional space proves superfluous.  Thus on all icons of the ``Protection of
the Virgin'' (Pokrov), the Mother of God is depicted directly in the space of the church.
Although she appeared invisibly to the people filling the church, here it sufficed that
she was seen by Andrei the Blessed and Epiphanius.  Similarly, in the very widespread
composition ``Apparition of the Mother of God and of the Apostles Peter and John to
Sergius of Radonezh,'' both the Mother of God and the apostles and Sergius with Mikhei
are depicted in one space.  Even on the ``Dormition'' icons described above, in cases
where they show the angel cutting off the hands of the impious Avfoniy (Ill.~69), the
angel is depicted in the same space as the apostles, since in the iconographer's
conception he acted in reality, in the real space and visibly to all.  The number of
such examples could be multiplied without difficulty.

\srcpage{155}

Similarly, there is no need to introduce two colours for distinguishing two spaces when
only the mystical space is depicted---for example, on icons of the ``Harrowing of Hell''
(Descent into Hades), in depictions of the righteous in paradise, and so on.

\figblock{70}{``Anthony of Rome.''  Icon from the Volga region.
  Sixteenth century.  State Tretyakov Gallery (GTG).}

The ``Transfiguration'' compositions deserve separate mention.  In them Christ and Moses
and Elijah conversing with him are in the same space as the Apostles Peter, John, and
James, who see and hear them.  Therefore the radiance depicted around Christ---often
including areas of blue sky with stars---cannot be regarded as an instance of the method
of sections.  This artistic device emphasises the divine nature of Christ and has no
bearing on the geometry of space.  Perhaps the blue starry sky, so reminiscent of the
blue of the mystical sky, was meant to symbolise Christ's connection with it.

The art of the Renaissance, which set as its principal aim the showing on the canvas of
the world surrounding man in its natural appearance---to show it so that the picture would
be as it were a ``window onto the material world''---was obliged to abandon the methods
characteristic of medieval art.  Among much else, Renaissance art abandoned drawing
methods.  In particular, the simultaneous depiction of two coexisting regions of
four-dimensional space by the method of sections was abandoned.  At the same time, the
surviving religious subject-matter remained the thematic basis of pictures.  The boundary
between the visible and invisible worlds came to be shown by the depiction of clouds.
Probably from this time there became established the notion that God, angels, and saints
``live on clouds''---that absurd notion which modern caricaturists entirely legitimately
exploit.  Geometrically it is obvious that real, existing clouds can only divide two
regions of real three-dimensional space and cannot separate two ``layers'' of a
four-dimensional world.  But the demand for ``naturalness'' led to an absurdity: the
``beyond-the-clouds'' world came to be depicted in the same colours as the earthly world,
and on the canvas ceased to differ from it.

\srcpage{156}

Thus, from the standpoint of geometric logic, artists who separated the real world from
the mystical by clouds were trying to drag a naive ``realism'' into a place where it is
out of place; this was a step backward compared to the logically impeccable creations of
the time of Theophanes the Greek, Rublev, and Dionysius.  Of course, we speak here only
of geometric logic and not of the artistic merits of the great creations of the masters
of the Renaissance.  These masters were pioneering new paths for art, making it ever more
tangible, one where even the deeply symbolic ``needed credibility, corporeality, the
three-dimensional quality of the image.''\footnote{L.~M.~Batkin.  ``Dialogichnost'
  italyanskogo Vozrozhdeniya'' [The Dialogic Character of the Italian Renaissance].---
  \emph{Sovetskoe iskusstvoznanie}, vyp.~2.  Moscow, 1977, p.~87.}

In the sixteenth and seventeenth centuries this method of separating real from mystical
space gradually passed into Russian painting as well, diminishing---together with other
causes---its philosophical depth.

\srcpage{157}

\figblock{71}{Nikifor Savin.  ``John the Baptist in the Wilderness.''
  Icon of the seventeenth century.  State Tretyakov Gallery (GTG).}

On the icon by Nikifor Savin (seventeenth century) ``John the Baptist in the Wilderness,''
Christ surrounded by angels is shown in the upper left corner of the icon, separated from
John only by billowing clouds (Ill.~71).  The space in which Christ stands differs from
the space in which John stands neither in colour nor in any other way.  The only hint
that Christ with the angels does not simply ``dwell'' on a cloud in real space is a
barely visible depiction of a certain firmament (?), on which the throne of Christ is
set, and this firmament is shown in the same colour as the earth beneath John's feet.  It
is perfectly clear that this contrived verisimilitude bears no comparison whatsoever with
the impeccable logic of the ``Dormition'' icons created by the great masters of ancient
painting.  The considerations presented here do not touch on the artistic merits of the
icon.

\srcpage{158}

In conclusion of the present section, let us adduce some further remarks on the depiction
of clouds in Old Russian painting.  In certain cases clouds played an essential role in
events occurring in real space, and then they were depicted alongside other real objects.
Thus, on the ``Dormition'' icon from Belozersky Monastery already mentioned (Ill.~69),
angels are shown delivering the apostles to the bier of the Mother of God, using clouds
for this purpose.\footnote{The motif of clouds swiftly carrying characters through the
  air frequently occurs in folklore, and such magical displacements are understood as
  occurring in real space.}  Both the clouds with the apostles and the angels bearing
them move in ordinary space; therefore the angels are shown as beings just as real as the
apostles.  They differ sharply from the angels forming part of the heavenly hosts
surrounding Christ, or the angels awaiting the Mother of God at the gates of paradise,
who are barely distinguishable and seem to merge with the background of the mystical
space.  Clouds differ from the mystical space in that instead of the geometrically precise
boundaries of the latter they have an indefinite ``billowing'' form; sometimes they are
sharply separated from the mystical space by colour\footnote{In this respect the most
  consistent is the icon ``Dormition'' of the fifteenth century from the State Tretyakov
  Gallery (see: V.~I.~Antonova, N.~Ye.~Mneva.  \emph{Katalog drevnerusskoy zhivopisi}
  [Catalogue of Old Russian Painting], vol.~I.  Moscow, 1963, No.~201, ill.~141);
  L.~M.~Evseeva, I.~A.~Kochetkov, V.~N.~Sergeyev.  \emph{Zhivopis' drevney Tveri}
  [Painting of Ancient Tver].  Moscow, 1974, ill.~26.} or by a shade of
colour.\footnote{In the sense indicated, the frescoes of the Rostov Kremlin churches are
  interesting, where real and mystical space frequently ``exchange'' colours.  Thus in
  the festal tier of the altar screen of the Church of John the Theologian, the sky is
  rendered now in blue, now in an almost white colour.  In cases where a depiction of the
  mystical space is also necessary, it is likewise rendered in blue or white, but always
  such that if the ordinary sky is rendered blue, the mystical space is assigned a white
  colour, and vice versa.}

Sometimes a cloud was depicted to provide ``support'' for the feet of the Mother of God
on the ``Protection'' icons (Pokrov), but even in this case such a cloud was conceived
as being together with the Mother of God in the real space of the church.  In both cases
described, as in other analogous ones, real and mystical space are clearly divided, and
clouds are in no sense boundaries between the visible and the invisible world.

\srcpage{160}

The depiction of clouds as the boundary between two worlds is used in Old Russian painting
from the second half of the sixteenth century onward and especially in the seventeenth
century.  Initially they are shown quite modestly, as if merely veiling the boundary
between two worlds, which was previously given by a clear line---sometimes emphasised by
colour (for example, a red stripe) and sometimes by stars arranged along it.  At this
stage artists still understand that they are depicting two different worlds, and therefore
use a distinct colour for each.

As regards the boundary between both spaces, in the seventeenth century it is often still
depicted by a not very substantial layer of clouds, usually painted in a colour that
coincides with neither of the two spaces and moreover not necessarily white.  All this as
it were underscores their role as the boundary between two incompatible spaces.

Later, in the seventeenth to nineteenth centuries, one and the same colours come to be
used for depicting both the ordinary and the mystical world, and the logic of the
construction is violated.  The two spaces are by tradition divided by clouds, now painted
almost naturalistically.  The substantial layer of billowing clouds grows; they cease to
serve as the sign of a conventional boundary between two spaces and seem like the most
natural cumulus clouds.  Christ, angels, and other inhabitants of the mystical space are
disposed

\srcpage{161}

\figblock{72}{Unfolding of a die.}

directly on the clouds, without any depiction of firm ground beneath their feet.  The
understanding of the fundamental difference between ordinary and mystical space vanishes.
The disappearance in painting of a clear separation of these two spaces, the forgetting
of the most natural method of their simultaneous depiction---this is a sign of serious
changes in the figurative structure of art.  Even in cases where the artist conveys only
the mystical space, he begins to show it as almost indistinguishable from the everyday.
As an example one may cite the icon ``Harrowing of Hell'' in the festal tier of the
iconostasis of the Assumption Cathedral of Ryazan (eighteenth century), where instead of
the conventional depiction of the mystical space, in particular the entrance to Hades,
accepted in earlier times, ordinary earth overgrown with some sort of bushes is shown,
and those freed by Christ move out of a natural underground cave.

The process briefly described here characterises one facet of the inevitable losses
connected with the transition from the depiction of the logical essence of a phenomenon
to the illusionistic aspiration to paint everything so that what is painted approximates
as nearly as possible the visible picture of the real world.

\noindent\textbf{3.}\quad\textbf{Unfoldings.}
In technical drawing, an \emph{unfolding} is a drawing of the cut-out pattern for sheet
material (for example, a sheet of iron).  Such a pattern is needed so that, having cut
the sheet material according to it and then processing it further---in particular by
bending---one obtains the required part.  An unfolding is therefore not a drawing of the
part itself but only a drawing of the blank needed for its manufacture.  Nonetheless,
such a drawing can give a very complete impression of the part for which it has been made.
Thus in the unfolding of an ordinary die, fold lines are drawn straight, so that the
material may be bent along them at a right angle to obtain the required solid figure (in
this case a cube) (Ill.~72).  Such a drawing gives absolutely complete information about
the markings on all six faces of the cube, which is impossible in an ordinary depiction
of a die, even in the three projections of an ordinary engineering drawing.

The aspiration of medieval artists to convey by means of the image they created as much
information as possible, to reveal in an object the qualities needed for understanding
the subject, was bound to direct their attention to the possibilities offered by methods
analogous to the drawing of unfoldings.

\srcpage{162}

A related device appears in miniatures principally in the depiction of buildings.  The
Middle Ages did not know drawing from life, and it is understandable that when reproducing
the appearance of a building from memory the artist did not feel bound to a single
viewpoint, but considered it his duty to give the most complete impression of it.

\figblock{73}{Engineering drawing of the south facade of the Assumption Cathedral of the
  Kremlin, Moscow.}

In Old Russian pictorial art, the method analogous to the drawing of unfoldings lacks
the strict geometric completeness of modern technical drawing, and its purpose is entirely
different.

In the sixteenth-century icon ``Anthony of Rome'' (Ill.~70), the church shown in the
background is depicted with obvious violations of the rules of pictorial rendering (though
not of draughtsmanship) familiar to us.  The central volume of the building is shown as
seen from the south; in particular, its lower part reveals the south portal.  The
bell-tower attached to it is ``unfolded,'' as a result of which its west side with the
west entrance to the church is visible.  Thus both the south side of the church and the
west side are shown in one plane; as is clear from what has been said, the iconographer
made use of the ``unfolding'' method here.

Every unfolding is by its nature a flat image and allows a building to be shown
``unfolded'' in the manner in which we have done with the die.  The artist's recourse to
this device leads to a strengthening of the planar character in the work---though this
circumstance was not the motivating reason for the medieval artist's

\srcpage{163}

recourse to the ``unfolding'' method.  Moreover, the artist shows what seems important to
him, omitting what is from his point of view inessential.

This device was employed fairly often.  In depicting the Assumption Cathedral of the
Moscow Kremlin, one artist stretched its five domes into a single line, whereas in fact
four of them stand at the corners of a square and the fifth is at its centre, at the
intersection of the diagonals (Ill.~67).\footnote{Deviations from the usual arrangement
  of domes are so rare that they can be disregarded (for example, the domes of the
  Assumption Church of Pskov-Pechery Monastery are arranged in a line because of the
  configuration of the underground caves beneath it).}  The master preferred to show the
number of domes, assuming that every viewer knows their mutual arrangement from his own
life experience.  Thus the viewer gained an impression of the exterior of the Assumption
Cathedral by combining the information received through contemplating the image (the
number of domes) with the a~priori information he possessed independently of the image,
from life experience---the type of configuration of the dome arrangement, given that their
number is known.

On the basis of the modern engineering drawing of the south facade of the same Assumption
Cathedral (Ill.~73), although it shows only three domes, the viewer, invoking a~priori
information (based on life experience), may conclude that the cathedral has five domes,
since three domes have never been placed above the main volume of a church with the two
flanking ones south of the plane of symmetry and parallel to the east--west line.  Clearly
the viewer sees the south facade, since the apses of the cathedral are shown on the right
in the drawing.  Thus both the conventional projection of the exterior in Ill.~67 and
the modern projection in Ill.~73 convey an impression of the cathedral's exterior only
by invoking additional information, which above was called a~priori.

If one raises the question of which image---medieval or modern---is the ``more correct,''
one must first agree on the criteria of correctness.  Let such a criterion be the
possibility of describing the exterior of the cathedral in words.  This criterion is
entirely reasonable for the Middle Ages, when even the builder often received from the
patron only a verbal description of the intended building.

\srcpage{164}

On the seventeenth-century icon (Ill.~67) both the north and south sides of the cathedral
are shown simultaneously.  This is evident from the fact that in its lower part people
move toward the baldachin from two opposite portals, which is only possible if they are
the north and south portals.  The complete symmetry of the image indicates the identical
appearance of the north and south facades.  Both walls terminate in four \emph{zakomary}
(semicircular gables), and each semicircle of a zakomar corresponds to one window of the
upper tier.  It should be noted in particular that the icon shows not eight zakomary but
only seven.  This is connected with the fact that the central dome of the cathedral
corresponds to two zakomary simultaneously---the northern and the southern.  Looking at
the lower part of the icon, one can see that the windows of the upper tier correspond to
those of the lower.  To obtain a more complete impression of the cathedral, one must
remember the actual arrangement of the five domes in space.

The modern engineering drawing of the south facade of the Assumption Cathedral shows only
the south wall with four zakomary, two rows of windows (four in each row), and one portal.
On the drawing the cathedral is crowned by three domes.  A more complete impression of
the cathedral requires additional information---namely, that the cathedral has five domes;
one must further take into account the usual character of the arrangement of these domes
in space and assume that the north and south sides of the cathedral are identical.

The comparison carried out shows that, from the point of view of the medieval viewer, the
modern engineering drawing has no particular advantages over the icon being considered.
If one confines oneself to considerations of the quantity of information, a depiction of
the cathedral in perspective---for example, in axonometric---could carry even more of it.
Clearly, the medieval artist in the case examined preferred the ``unfolding'' not because
of the quantity of information it carries but because of the qualitative aspect of the
information, its orientation; for no other type of image is capable of conveying what the
artist considered most important: to show the south and north sides of the cathedral
simultaneously, the exterior and interior simultaneously.  This is possible only by
invoking drawing methods---the cross-section and the unfolding.  The importance of just
this compositional solution follows from the content of the icon.  The artist's aim was
to depict the solemn act in the Moscow Assumption Cathedral, not in any other church or
room; therefore it was necessary to show the external features of the Assumption Cathedral.
The comparison of icon and modern engineering drawing cited above convinces one that this
was done with sufficient clarity, additionally underlined by the placement of the icon of
the ``Dormition'' above the baldachin.  The north and south portals had to be shown
simultaneously and in the same way, to emphasise that the secular authorities (the Tsar
and the boyars accompanying him) coming from the south, and the ecclesiastical authorities
(the Patriarch with the clergy) coming from the north, are equally significant.  Beyond
what has been said, the method of depicting the building chosen by the artist permitted
the whole composition to be given a symmetrical character, which in itself may have
appeared to the artist sufficiently valuable.

Thus the use of ``unfoldings'' in Old Russian painting may be dictated by different
reasons---both informational considerations and compositional requirements.

\srcpage{165}

The Old Russian artist had a wide choice of methods for conveying the geometry of the
external world: he had the right not only to displace the inevitable distortions involved
in the perspectival rendering of the external world in the direction he desired, but, if
this seemed expedient, to abandon altogether any system of perspective and turn to drawing
methods.

\figblock{74}{``Miracle of Theodore Tiron.''  End of the fifteenth century.
  State Russian Museum (GRM).}

A vivid example of the fruitfulness of such a synthetic method is provided by the icon
``Miracle of Theodore Tiron'' (Ill.~74), in which three drawing devices are simultaneously
used: the upper part of the well into which Theodore descended is shown in a conventional,
right-angle-rotated position; in the lower part of the icon his adventures underground
are shown by means of a cross-section; and in the upper left corner the mystical sky is
depicted with a blessing hand of the deity (that is, the method of sectioning
four-dimensional space is employed).  One may safely say that the drawing methods bear
the whole composition of this icon, which was required simultaneously to show the mystical
sky, the earth, and the underworld.

\srcpage{166}

In the example mentioned here, the use of drawing methods had as its principal aim the
expansion of the artist's possibilities in a narrative subject---in this case, to tell
the story of the descent of Theodore Tiron underground.  At times drawing methods may be
the only means of conveying some important idea.

As is well known, in medieval compositions of the ``Crucifixion'' the ground at the base
of the cross is shown in cross-section, with the skull and bones of Adam depicted in the
grave directly beneath the cross.  Such a simultaneous depiction of the remains of Adam
and of Christ crucified above his grave symbolises the cardinal idea of Christ's
redemption of Adam's sin.  In the Renaissance, when drawing methods were banished from
the artist's arsenal, the pictorial means for this most important Christian idea of the
sacrificial redemption had been lost.  The forced abandonment of medieval iconography was
probably felt by artists as a loss.  This led to the appearance of different types of
``Crucifixion'': alongside the most ``natural,'' in which Adam's skull does not figure at
all, there appeared ones where the skull lies at the base of the cross or some distance
from it on the surface of the earth, and even ones where many skulls and bones are
scattered on Golgotha not only at the foot of the crucified Christ but also at the crosses
of the thieves.  The striving for outward plausibility led to absurdities, since it is
impossible to explain convincingly the appearance of the remains of the patriarch Adam on
Golgotha.  The possibility of the direct conveyance of the idea of Christ's sacrificial
redemption disappeared.  The figurative devices no longer correspond to the inner meaning
of the symbolic idea.

\srcpage{167}

In the light of what has been said, it is clear that the condescending tone sometimes
encountered in the literature toward the medieval artist with his drawing methods is
entirely out of place.  Drawing methods were highly reasonable, were not connected with
any ``inability,'' and constituted a legitimate feature of pictorial art.
